\clearpage
\section{维修员松开前进键时后退}

按住 W 走几步，再松开。旧版中，蓝色维修员看起来退回了半步。第三人称相机本身也在跟随主角，因此需要把三个位置分开记录。游戏每帧都写下角色的位置、骨盆在世界中的位置和相机的位置。角色坐标反映物理移动，骨盆减去角色坐标反映动画让身体移动了多少。前进方向沿场景的 X 轴，这一段可以直接比较两帧的 X 分量。

物理角色在切到待机的这几帧里仍向前移动了约 4.6 毫米。模型骨盆却相对角色后退了 23.66 厘米，同时向侧面偏。检查模型后可以看出，两段动作把同一副身体放在了不同的地方。走路时固定名为 \code{root} 的运动根，避免动画与碰撞体重复前进；待机的水平偏移存在它的子骨骼 \code{hips} 里，固定根并不能消掉这一层位移。

重新导入时，用走路完整迈完左右两步的骨盆平均位置作为站立位置。导入器把每个原地片段的骨盆关键帧平移相同的水平距离，站立高度、手臂角度和脚步周期仍取自作者动作。引擎按动作表里的时间衔接两个姿势。新版实机松手后的腰部后退 0.59 厘米，计入侧向偏移后的总距离为 0.83 厘米。测量区间从切入待机前一帧开始，到待机开始 0.2 秒后的第一帧结束。它包含了两段动作各自的自然摆动。

\loc{scripts/retarget_station_motion.py --- align_standing_clips}
\loc{Redflag/src/Redflag/Animation/ModelAnimationPlayer.cpp --- PlayClip}

\textbf{长按后退碰到护栏。} 角色确实走出了约 6.26 米。走到护栏前停下是正常的，随后按 W 往反方向走却一直贴住原地。减速时产生的轻微重叠使下一帧的移动从墙面内部开始。旧查询把这种命中一律当成前面也被堵住。新版让 PhysX 返回从墙里退出的方向。角色往外走时，查询会继续寻找前方的另一堵墙；往内走仍会被第一堵墙挡住。测试用两堵互相独立的墙验证两个方向，也复查相机原有的碰撞查询。

同一次游戏操作里，角色重新按 W 离开护栏后的位移从 0 增至 1.33 米。完整的旧版与新版录像和逐帧轨迹见 \path{output/validation/station-motion-stop/}。两次均以 Vulkan 在 1280 × 720 画面运行；这是动作和位置验证。

\loc{Redflag/src/Redflag/Physics/PhysicsWorld.cpp --- PhysicsWorld::SweepSphere}
\loc{tests/TerrainTraversalTests.h --- CheckTraversal}
